\documentclass[a4paper]{article}     

\usepackage{amsfonts}
\usepackage{amsmath}
\usepackage{amssymb}
\usepackage{graphicx}
\usepackage{color}

\newcommand{\R}{\mathbb{R}}
\newcommand{\llim}{\lim\limits}

\begin{document}

\noindent \large Auteur: Sadia Nazary \\
\noindent \large Studierichting: Informatica\\
\noindent \large Oefeningenreeks 7A nr. 18.1\\

\medskip

\normalsize

$f:\R^3\longrightarrow\R:(x,y,z)\mapsto f(x,y,z)=
\left\{
\begin{array}{lll}
\dfrac{xy+yz+zx}{x^2+y^2+z^2} & \text{als} & (x,y,z)\neq(0,0,0)\\
0 & \text{als} & (x,y,z)=(0,0,0)
\end{array}
\right .$\\

We onderzoeken de functie in $O=(0,0,0)$.\\

\medskip

Continuiteit in $O$:\\

Neem de rechte $x=y=z=t$.\\

Dan geldt voor $t\neq0$:\\

$f(t,t,t)
=
\dfrac{t^2+t^2+t^2}{t^2+t^2+t^2}$\\

$=
\dfrac{3t^2}{3t^2}$\\

$=
1$\\

Dus:\\

$\llim_{t\rightarrow0}f(t,t,t)=1$\\

Maar:\\

$f(0,0,0)=0$\\

Bijgevolg is $f$ niet continu in $O$.\\

\medskip

Partiele afgeleiden in $O$:\\

$\dfrac{\partial f}{\partial x}(0,0,0)
=
\llim_{h\rightarrow0}
\dfrac{f(h,0,0)-f(0,0,0)}{h}$\\

$=
\llim_{h\rightarrow0}
\dfrac{0-0}{h}$\\

$=
0$\\

Op dezelfde manier geldt:\\

$\dfrac{\partial f}{\partial y}(0,0,0)=0$\\

$\dfrac{\partial f}{\partial z}(0,0,0)=0$\\

Dus $f$ is partieel afleidbaar in $O$.\\

Omdat $f$ niet continu is in $O$, is $f$ niet continu partieel afleidbaar in $O$.\\

\medskip

Afleidbaarheid ten opzichte van een willekeurige richting:\\

Neem een willekeurige richting $(a,b,c)\neq(0,0,0)$.\\

Dan geldt voor $t\neq0$:\\

$f(ta,tb,tc)
=
\dfrac{t^2ab+t^2bc+t^2ca}{t^2a^2+t^2b^2+t^2c^2}$\\

$=
\dfrac{ab+bc+ca}{a^2+b^2+c^2}$\\

Dus:\\

$\llim_{t\rightarrow0}
\dfrac{f(ta,tb,tc)-f(0,0,0)}{t}$\\

$=
\llim_{t\rightarrow0}
\dfrac{1}{t}
\dfrac{ab+bc+ca}{a^2+b^2+c^2}$\\

Deze limiet bestaat in het algemeen niet.\\

Dus $f$ is niet afleidbaar ten opzichte van een willekeurige richting in $O$.\\

\medskip

Differentieerbaarheid in $O$:\\

Elke differentieerbare functie is continu.\\

Omdat $f$ niet continu is in $O$, is $f$ niet differentieerbaar in $O$.\\

\medskip

Besluit:\\

$f$ is niet continu in $O$.\\

$f$ is partieel afleidbaar in $O$.\\

$f$ is niet continu partieel afleidbaar in $O$.\\

$f$ is niet afleidbaar ten opzichte van een willekeurige richting in $O$.\\

$f$ is niet differentieerbaar in $O$.\\

\begin{thebibliography}{999}
%\bibitem{a} Referentie
\end{thebibliography}

\end{document}
